\section{Sharp list bounds for polynomial Riccati solutions}
\label{app:riccati-cross-ratio}

The classical constant-cross-ratio property of Riccati equations gives
an ordinary-list bound even when all agreement coordinates are singular.
For a challenge-independent nonlinear equation it also gives a linear
full-support exceptional-label bound. The parametrization and its gcd
form are classical; see Gasull, Torregrosa, and Zhang
\cite[Lemmas 6--7]{gasull-riccati}. We prove the finite-characteristic
statements and the agreement consequences used here.

\begin{theorem}[Riccati lists and solution counts]
\label{thm:riccati-cross-ratio}
Let $D\ge1$, and let $E$ have characteristic zero or $p>D$.
Consider one fixed equation
\begin{equation}
 a(X)P'=b_0(X)+b_1(X)P+b_2(X)P^2,\qquad a\ne0,
 \label{eq:fixed-riccati}
\end{equation}
with coefficients in $E[X]$, with no restriction on their degrees.
On any $n$ distinct evaluation points, the number of distinct solutions
of degree at most $D$ agreeing with an arbitrary received word in at
least $A>D$ coordinates is at most
\begin{equation}
 \left\lfloor\frac{n}{A-D}\right\rfloor.
 \label{eq:riccati-list}
\end{equation}
This bound is attained over prime fields at unbounded lengths and
positive limiting rate and agreement gap. If $b_2\ne0$, the total
number of degree-at-most-$D$ polynomial solutions is at most
\begin{equation}\label{eq:riccati-total-count}
 C_D=\begin{cases}
 D+2,&\text{in characteristic zero or }p>D+1,\\
 2D+2,&\text{if }p=D+1.
 \end{cases}
\end{equation}
The bound $D+2$ is sharp. The boundary bound $2D+2$ is attained
at $p=2,D=1$.
\end{theorem}
\begin{proof}
\emph{Total candidate count.}
Assume $b_2\ne0$. Work at a transcendental Taylor center $t$, where
$a(t)b_2(t)\ne0$, and write the equation as
$P'=F(X,P)$, with $F=\alpha+\beta u+\gamma u^2$ and $\gamma\ne0$.
Put $\Delta=\partial_X+F\partial_u$ and $F_j=\Delta^j u$.
For $j<p$, or in characteristic zero, $F_j$ has $u$-degree $j+1$
and leading coefficient $j!\gamma^j$. This follows inductively:
$F\partial_u$ gives the highest degree and $\partial_X$ gives a
lower degree. In characteristic zero or $p>D+1$, every solution
satisfies
\[
 F_{D+1}(t,P(t))=P^{(D+1)}(t)=0.
\]
This is a nonzero polynomial of degree $D+2$ in $P(t)$.
Distinct polynomial solutions have distinct values at the
transcendental $t$, proving the first count.

If $p=D+1$, use instead the truncated reconstruction
\[
 T(U,X)=\sum_{j=0}^D F_j(t,U)\frac{(X-t)^j}{j!}.
\]
Every candidate is $T(P(t),X)$. Its $U$-degree is $D+1$, with
leading coefficient $\gamma(t)^D(X-t)^D$. The residual
$a\partial_XT-b_0-b_1T-b_2T^2$ therefore has $U$-degree exactly
$2D+2$, with nonzero leading coefficient
$-b_2(X)\gamma(t)^{2D}(X-t)^{2D}$. Any nonzero $X$-coefficient
of this residual is a nonzero polynomial in $U$ of degree at most
$2D+2$ vanishing at every candidate value $P(t)$. This proves the
boundary count. In particular, in positive characteristic the total
count is at most $2D+2\le2p$.

\emph{Lists by intersection multiplicities.}
For $M\ge2$ distinct candidates and a domain coordinate $x$, put
\[
 C_x=\sum_{i<j}\operatorname{ord}_x(P_i-P_j).
\]
The root bound gives $\sum_xC_x\le D\binom M2$. If $h$ candidates
match a given received value at $x$, we claim
\begin{equation}\label{eq:riccati-local-multiplicity}
 (h-1)(M-1)\le2C_x.
\end{equation}
If $b_2=0$, differences solve one homogeneous linear equation.
The ratio of two nonzero differences has derivative zero and
rational degree at most $D<p$, and hence is constant. The same is
true in characteristic zero. Thus the candidates have form $P_0+cW$,
and their values at $x$ are either all distinct or all equal;
\eqref{eq:riccati-local-multiplicity} follows immediately.

Now suppose $b_2\ne0$. Differences satisfy
\[
 a(P_i-P_j)'=(b_1+b_2(P_i+P_j))(P_i-P_j),
\]
so every cross ratio has derivative zero. In characteristic zero
it is constant and different from $0,1,\infty$. A colliding pair
and two candidates outside its value class would force that cross
ratio to be one. Therefore either all values are distinct or one
value occurs at least $M-1$ times. For
$h\in\{0,1,M-1,M\}$, the inequality follows from
$C_x\ge\binom h2$.

In positive characteristic a cross ratio need not be constant.
However, any zero of a nonzero rational function with derivative
zero has order divisible by $p$. Choose a colliding pair and two
candidates outside its value class. The denominator of their
cross-ratio identity is nonzero at $x$, and the numerator of the
cross ratio minus one, up to sign, is the product of the two
within-pair differences. If the two outside values were distinct,
this product would have positive order at most $D<p$, impossible.
Consequently every collision pattern other than those already
handled has exactly two value classes, both of size at least two.

Let the matching class have size $h$ and the other class have size
$b=M-h$. Let $\mu,\nu$ be the minimum positive vanishing orders of
differences within these classes. The same cross-ratio argument
for every choice of two within-class pairs gives $\mu+\nu\ge p$.
Write $A_0=\binom h2$, $B_0=\binom b2$, so
$C_x\ge A_0\mu+B_0\nu$. If $h\le b$, already
\[
 2(A_0+B_0)-(M-1)(h-1)=(b-1)(b-h+1)\ge0.
\]
If $h\ge b$, minimizing with $\mu,\nu\ge1$ and
$\mu+\nu\ge p$ gives $C_x\ge A_0+(p-1)B_0$, and hence
\[
 2C_x-(M-1)(h-1)\ge(b-1)(pb-M+1)\ge0,
\]
because $b\ge2$ and the total-count bound gives $M\le2p$.
An unmatched received value has $h=0$. This proves
\eqref{eq:riccati-local-multiplicity} in all cases.
Summing it gives $AM\le n+DM$, proving
\eqref{eq:riccati-list}. The case $M=1$ is immediate, and the same
proof allows real agreement thresholds.

For sharpness of $D+2$, choose $a$ split and squarefree of degree
$D+1$. The equation $aP'=a'P-P^2$ has the $D+2$ solutions $0$ and
$a/(X-r)$, one for each root $r$ of $a$. At the boundary $p=2,D=1$,
$(X^2+X)P'=P^2+P$ has exactly $0,1,X,X+1$ as its four solutions.
\end{proof}

\begin{proposition}[Sharp prime-field lists]
\label{prop:riccati-sharp-list}
Fix $M\ge3$ and a prime $p\equiv1\pmod{2M}$. Put
\[
 s=\frac{p-1}{2M},\quad n=p-1,\quad D=(M-2)s,\quad A=Ms.
\]
There is a nonlinear equation~\eqref{eq:fixed-riccati} over $\F_p$
and a word on $\F_p^*$ whose list of degree-at-most-$D$ solutions
at agreement threshold $A$ has exactly $M=n/(A-D)$ elements.
\end{proposition}
\begin{proof}
Partition the image of $X^s$ on $\F_p^*$ into
$a_1,\ldots,a_M,b_1,\ldots,b_M$ and set
\[
 R(Y)=\prod_{i=1}^M(Y-a_i),\quad T(Y)=Y^{M-1},\quad
 C_i(Y)=\frac{R(Y)}{Y-a_i}-T(Y),\quad P_i(X)=C_i(X^s).
\]
Leading terms cancel, so $\deg P_i\le D$. On the fiber over $a_j$
choose received value $-T(a_j)$, and on the fiber over $b_j$ choose
$C_j(b_j)$. Exactly $M-1$ candidates match a root fiber, since the
remaining difference is $R'(a_j)\ne0$. The values $C_i(b_j)$ are all
distinct, since $R(b_j)\ne0$ and the denominators $b_j-a_i$ are
distinct. Each $P_i$ therefore agrees in exactly $(M-1)s+s=A$
coordinates, and the polynomials are distinct.

The polynomials $U_i=R/(Y-a_i)$ satisfy $RU_i'=R'U_i-U_i^2$.
Translating $U=C+T$ gives
\[
 RC'=(R'T-T^2-RT')+(R'-2T)C-C^2.
\]
Substitute $Y=X^s$ and multiply the right side by $sX^{s-1}$.
This is one polynomial Riccati equation for all the $P_i$, with
nonzero quadratic coefficient. Since $p>2D$, the preceding theorem
shows that no additional solution can enter this list. Primes in
the progression $1\pmod{2M}$ give unbounded lengths for fixed $M$.
The code rate $(D+1)/n$ tends to $(M-2)/(2M)$ and the capacity gap
$(A-D-1)/n$ tends to $1/M$.
\end{proof}

\begin{corollary}[Fixed-equation full-support exceptional labels]
\label{cor:riccati-fixed-mca}
Under Theorem~\ref{thm:riccati-cross-ratio}, assume $b_2\ne0$.
For a received line $f+zg$ on $n$ distinct points and $A>D$, the
number of labels possessing a solution of the fixed equation with
at least $A$ agreements and bad full agreement support is at most
\begin{equation}
 \frac{2n(n+C_D)}{A-D}.
 \label{eq:riccati-fixed-mca}
\end{equation}
Here a support is bad if no polynomial of degree at most $D$
interpolates $g$ on the entire support. Thus the bound is $O(n)$
when $A-D$ is a fixed positive fraction of $n$.
\end{corollary}
\begin{proof}
Put $\delta=A-D$. A candidate's persistent coordinates have
$g_i=0$ and $P(x_i)=f_i$. Call a candidate heavy if it has at least
$D+\delta/2$ persistent coordinates. These candidates form a list
for $f$ at that threshold, so there are at most $2n/\delta$ by
\eqref{eq:riccati-list}. Each has at most $n$ bad labels: a bad
support must contain an agreement at a coordinate with $g_i\ne0$,
and this coordinate determines a unique label. Otherwise $g=0$
on the support, contradicting badness.

A light candidate needs more than $\delta/2$ nonpersistent
agreements at every nearby label. It has at most $n$ such
coordinate-label incidences altogether, so it has at most
$2n/\delta$ nearby labels. There are at most $C_D$ candidates in
total, with $C_D$ as in~\eqref{eq:riccati-total-count}. Adding the heavy and light bounds proves the claim.
\end{proof}

\paragraph{Characteristic guard and scope.}
The constant-cross-ratio argument alone cannot justify the weaker
guard $p>D$; intersection multiplicities are essential. In characteristic three,
put $D=2$, $U=1+X+X^2$, $V=X+X^2$, $W=1+X$, and $N=UV$.
The equation with $a=N$, $b_0=0$, $b_1=N'$, and
$b_2=U'V-UV'$ has the solutions $0,U,V,W$. Their values at zero
are $0,1,0,1$, contradicting the collision property. Their cross
ratio is a nonconstant rational function with zero derivative.
Exact enumeration gives five degree-at-most-two solutions, also
exceeding $D+2$. This is consistent with the boundary value
$C_D=2D+2$ in~\eqref{eq:riccati-total-count}.

The ordinary-list bound applies separately at each label of a
challenge-dependent Riccati equation. The full-support bound
\eqref{eq:riccati-fixed-mca} requires a single fixed equation and
does not extend by this argument to challenge-dependent families.
The checks in \path{research/riccati_cross_ratio/} enumerate
$183{,}581$ candidate polynomials, test $616$ received lines, and
verify nine prime-field equality fixtures for
Proposition~\ref{prop:riccati-sharp-list}, together with the
characteristic-three negative control. A further boundary-characteristic
checker enumerates $1{,}084{,}807$ polynomials for $549$ equations,
checks $2{,}790$ coordinate-family instances of the multiplicity
inequality, and includes nonconstant cross ratios and the sharp
characteristic-two boundary count. An independent extension-field
implementation exhausts all $729$ degree-at-most-two polynomials
over $\mathbb F_9$ and checks $85$ received lines, including bad
full-support labels at $p=D+1$. These restricted results
do not settle the general first-order conjecture.
